% Generated draft: requires mathematical and source review.
\documentclass[11pt]{article}
\usepackage[margin=25mm]{geometry}
\usepackage{fontspec}
\setmainfont{Latin Modern Roman}
\usepackage{amsmath,amssymb,mathtools,unicode-math}
\setmathfont{Latin Modern Math}
\usepackage{enumitem,needspace,xurl,hyperref}
\setlength{\parindent}{0pt}
\setlength{\parskip}{5pt}
\setlength{\emergencystretch}{3em}
\newcommand{\N}{\mathbb{N}}
\newcommand{\Z}{\mathbb{Z}}
\newcommand{\Q}{\mathbb{Q}}
\newcommand{\R}{\mathbb{R}}
\newcommand{\C}{\mathbb{C}}
\newcommand{\F}{\mathbb{F}}
\newcommand{\A}{\mathbb{A}}
\newcommand{\PP}{\mathbb{P}}
\newcommand{\E}{\mathbb{E}}
\newcommand{\eps}{\varepsilon}
\newcommand{\dd}{\,\mathrm{d}}
\begin{document}
\section*{UnsolvedMath ID: 701286}
\textbf{Source problem code:} catalogue entry
\subsection{Definitions and mathematical statement}
Let \(F\) be a field and let \(n\geq 1\) be an integer. Write \(\operatorname{GL}(n,F)\) for the group of invertible \(n\times n\) matrices over \(F\). A permutation matrix is a matrix \(P\) for which there is a permutation \(\sigma\) of \(\{1,\ldots,n\}\) such that \(P_{ij}=1\) when \(i=\sigma(j)\), and \(P_{ij}=0\) otherwise. For an \(n\times n\) matrix \(B\) over \(F\), its characteristic polynomial is \(\chi_B(t)=\det(tI_n-B)\), and its minimal polynomial \(m_B(t)\) is the monic polynomial of least positive degree satisfying \(m_B(B)=0\). Call \(B\) cyclic if \(m_B(t)=\chi_B(t)\).

\textbf{Conjecture (J. G. Thompson, Kourovka Notebook 16.95).} For every field \(F\), every integer \(n\geq 1\), and every \(A\in\operatorname{GL}(n,F)\), there exists an \(n\times n\) permutation matrix \(P\) such that
\[
 m_{AP}(t)=\chi_{AP}(t)\qquad\text{in }F[t].
\]
The supplied formulation places no restriction on the characteristic or cardinality of \(F\).
\subsection{Short English statement}
The conjecture asserts that every invertible square matrix over any field can have its columns permuted so that its minimal and characteristic polynomials coincide.
\subsection{Sources}
\begin{itemize}
\item Local UnsolvedMath registry, permanent ID \(701286\), source code \texttt{catalogue entry}, title \emph{Kourovka Notebook Problem 16.95}. The supplied record labels the problem open; this status has not been independently verified.
\item \emph{The Kourovka Notebook (Unsolved Problems in Group Theory)}, problem 16.95, proposed by J. G. Thompson in the 16th issue (2006). The supplied background also identifies the 21st issue edition (2026), \texttt{arXiv:1401.0300v48}.
\item The supplied collection description attributes the editing of the Kourovka Notebook, Issues 1--21 (1965--2026), to E. I. Khukhro and V. D. Mazurov.
\end{itemize}
\end{document}
